\documentclass{article}

% if you need to pass options to natbib, use, e.g.:
%     \PassOptionsToPackage{numbers, compress}{natbib}
% before loading neurips_2026

% The authors should use one of these tracks.
% Before accepting by the NeurIPS conference, select one of the options below.
% 0. "default" for submission
\usepackage{neurips_2026}
% the "default" option is equal to the "main" option, which is used for the Main Track with double-blind reviewing.
% 1. "main" option is used for the Main Track
%  \usepackage[main]{neurips_2026}
% 2. "position" option is used for the Position Paper Track
%  \usepackage[position]{neurips_2026}
% 3. "eandd" option is used for the Evaluations & Datasets Track
 % \usepackage[eandd]{neurips_2026}
 % if you need to opt-in for a single-blind submission in the E&D track:
 %\usepackage[eandd, nonanonymous]{neurips_2026}
% 4. "creativeai" option is used for the Creative AI Track
%  \usepackage[creativeai]{neurips_2026}
% 5. "sglblindworkshop" option is used for the Workshop with single-blind reviewing
 % \usepackage[sglblindworkshop]{neurips_2026}
% 6. "dblblindworkshop" option is used for the Workshop with double-blind reviewing
%  \usepackage[dblblindworkshop]{neurips_2026}

% After being accepted, the authors should add "final" behind the track to compile a camera-ready version.
% 1. Main Track
 % \usepackage[main, final]{neurips_2026}
% 2. Position Paper Track
%  \usepackage[position, final]{neurips_2026}
% 3. Evaluations & Datasets Track
 % \usepackage[eandd, final]{neurips_2026}
% 4. Creative AI Track
%  \usepackage[creativeai, final]{neurips_2026}
% 5. Workshop with single-blind reviewing
%  \usepackage[sglblindworkshop, final]{neurips_2026}
% 6. Workshop with double-blind reviewing
%  \usepackage[dblblindworkshop, final]{neurips_2026}
% Note. For the workshop paper template, both \title{} and \workshoptitle{} are required, with the former indicating the paper title shown in the title and the latter indicating the workshop title displayed in the footnote.
% For workshops (5., 6.), the authors should add the name of the workshop, "\workshoptitle" command is used to set the workshop title.
% \workshoptitle{WORKSHOP TITLE}

% "preprint" option is used for arXiv or other preprint submissions
 % \usepackage[preprint]{neurips_2026}

% to avoid loading the natbib package, add option nonatbib:
%    \usepackage[nonatbib]{neurips_2026}

\usepackage[utf8]{inputenc} % allow utf-8 input
\usepackage[T1]{fontenc}    % use 8-bit T1 fonts
\usepackage{hyperref}       % hyperlinks
\usepackage{url}            % simple URL typesetting
\usepackage{booktabs}       % professional-quality tables
\usepackage{amsfonts}       % blackboard math symbols
\usepackage{amsmath,amssymb}
\usepackage{nicefrac}       % compact symbols for 1/2, etc.
\usepackage{microtype}      % microtypography
\usepackage{xcolor}         % colors
\usepackage{graphicx}
\usepackage{multirow}
\usepackage{array}
\usepackage{siunitx}
\usepackage{caption}
\usepackage{subcaption}
\usepackage{listings}
\usepackage{algorithm}
\usepackage{algpseudocode}

\lstdefinestyle{cuda}{
  basicstyle=\ttfamily\footnotesize,
  keywordstyle=\color{blue!70!black}\bfseries,
  commentstyle=\color{green!50!black}\itshape,
  stringstyle=\color{red!70!black},
  showstringspaces=false,
  breaklines=true,
  numbers=left,
  numberstyle=\tiny\color{gray},
  frame=single,
  language=C++,
}

\title{AccelEval: Evaluating and Decomposing LLM-Generated CUDA Code at Scale}


% The \author macro works with any number of authors. There are two commands
% used to separate the names and addresses of multiple authors: \And and \AND.
%
% Using \And between authors leaves it to LaTeX to determine where to break the
% lines. Using \AND forces a line break at that point. So, if LaTeX puts 3 of 4
% authors names on the first line, and the last on the second line, try using
% \AND instead of \And before the third author name.


\author{%
  Chen Dong \\
  % Affiliation \\
  \texttt{zijieteam@gmail.com} \\
  % \And
  % Co-Author \\
  % Affiliation \\
  % \texttt{email} \\
}


\begin{document}


\maketitle


\begin{abstract}
We introduce \textsc{AccelEval}, a benchmark for evaluating large language models (LLMs) on translating sequential CPU programs into optimized CUDA kernels. Unlike prior kernel-generation benchmarks that focus on isolated GPU primitives, \textsc{AccelEval} stresses end-to-end acceleration across \textbf{42 production-style workloads} spanning scientific computing, graph analytics, financial engineering, dynamic programming, and operations research. Each task provides a C reference implementation, a three-scale input generator, and a deterministic validation oracle. We report results for six state-of-the-art models under a zero-hint prompting regime (L3), comparing compilation success, numerical correctness, and end-to-end speedup over optimized CPU baselines. The strongest model (Gemini~3.1~Pro) achieves an 88.9\% pass rate with a median speedup of $316\times$, while other models cluster between 47\% and 78\%. The gap concentrates on \emph{correctness-critical} HPC tasks (sparse linear algebra, multigrid, molecular dynamics), suggesting LLMs still struggle with fine-grained numerical reasoning under parallel execution.
\end{abstract}



\section{Introduction}
\label{sec:intro}

GPU acceleration is the dominant path to scaling scientific and industrial workloads, yet authoring efficient CUDA remains the domain of specialists. Recent LLMs exhibit strong code-generation capability, raising the question: \emph{can an LLM automate the CPU-to-CUDA porting process end-to-end?}

Existing benchmarks answer this only partially. KernelBench~\citep{kernelbench} measures single-kernel generation; BabelTower~\citep{babeltower} targets code-level translation without correctness oracles. None evaluate \emph{realistic workloads} with full harnesses, multi-scale inputs, and performance comparison against tuned CPU baselines.

\paragraph{Contributions.}
\begin{enumerate}
  \item We release \textsc{AccelEval}, a curated suite of \textbf{42 CPU-to-CUDA tasks} across 20+ algorithmic categories, each with a reference implementation, binary I/O harness, validation oracle, and three input scales.
  \item We define a three-level prompting protocol (L1: full hints; L2: brief hints; L3: zero hints) that disentangles raw capability from prompt engineering.
  \item We present the first large-scale evaluation of six leading LLMs (Gemini~3.1~Pro, Claude~Opus~4.6, Qwen~3.6~Plus, Kimi~K2.5, GLM~5.1, DeepSeek~V3.2) under L3, reporting compile rate, pass rate, and end-to-end speedup.
  \item Beyond leaderboard scores, we introduce a \textbf{decomposition pipeline}: an auto-detected \emph{optimization pattern database} combined with \emph{pairwise diff attribution} over multi-turn agent runs, enabling causal claims of the form ``pattern $P$ contributes $\Delta$ speedup.''
\end{enumerate}

\section{Related Work}
\label{sec:related}

\paragraph{LLM code generation benchmarks.}
HumanEval~\citep{humaneval} and MBPP~\citep{mbpp} established text-to-code evaluation for Python. SWE-bench~\citep{swebench} pushed to repository-level fixes. None target GPU programming.

\paragraph{GPU code benchmarks.}
KernelBench~\citep{kernelbench} is the closest peer. It evaluates single-kernel replacements inside a PyTorch harness and reports speedups relative to \texttt{torch.compile}. Our setting differs in two ways: (i)~we evaluate \emph{whole-program} translation including I/O, memory management, and multi-kernel orchestration; (ii)~our reference is a pinned CPU implementation, giving a fixed absolute baseline.

\paragraph{Automatic parallelization and translation.}
Prior work on source-to-source compilation (OpenACC, Bones~\citep{bones}, Par4All~\citep{par4all}) relies on static analysis and pragma hints. LLM-based translators like BabelTower~\citep{babeltower} and CodeRosetta~\citep{coderosetta} show promise but lack end-to-end execution validation.

\section{Problem Formulation}
\label{sec:formulation}

We cast CPU-to-CUDA optimization as a \emph{strategy recommendation}
problem over a structured algorithm space. This framing lets us study
LLMs not only as black-box generators, but as policies operating in a
well-defined decision space, enabling both principled baselines and
decomposition of where their advantage comes from.

\subsection{Problem Space $\mathcal{P}$}
Let $\mathcal{P} = \{Q_1, \ldots, Q_n\}$ be a set of CPU reference
snippets (one per \textsc{AccelEval} task). An \emph{affinity function}
$\rho : \mathcal{P} \times \mathcal{P} \to \mathbb{R}_{\geq 0}$
measures structural similarity, making $(\mathcal{P}, \rho)$ a weighted
graph. In our instantiation, $\rho$ is induced by a kernel over
static-analysis features (Section~\ref{sec:feature}).

\subsection{Feature Space $\mathcal{F}$}
\label{sec:feature}
A feature extractor $f : \mathcal{P} \to \mathcal{F} \subseteq
\mathbb{R}^d$ maps each problem to a vector
$\mathbf{q} = f(Q)$ capturing loop-nest depth, stride-1 access ratio,
arithmetic intensity, indirection level, reduction presence, etc.
Affinity is realized via a kernel:
\begin{equation}
\rho(Q_i, Q_j) = \kappa_\mathcal{F}\bigl(f(Q_i), f(Q_j)\bigr).
\label{eq:affinity}
\end{equation}

\subsection{Algorithm Space $\mathcal{A}$}
Let $\mathcal{T} = \{t_1, \ldots, t_m\}$ be a catalog of
transformation rules (Table~\ref{tab:trans}). Each rule $t_i$ carries
an intensity level $s_i \in \{0, 1, \ldots, M_i\}$, with $s_i = 0$
meaning inactive. A \emph{strategy vector} is
\begin{equation}
\mathbf{s} = (s_1, \ldots, s_m) \in \mathcal{A}
   = \prod_{i=1}^{m}\{0, 1, \ldots, M_i\}.
\label{eq:strategy}
\end{equation}

\begin{table}[ht]
\centering
\small
\caption{Example transformation rules with intensity semantics.}
\label{tab:trans}
\begin{tabular}{@{}clll@{}}
\toprule
$s_i$ & \textbf{Rule} & \textbf{Intensity semantics} & $M_i$ \\
\midrule
$s_1$ & Shared-memory tiling & $0$: off; $\ell$: tile level $\ell$ & 4 \\
$s_2$ & Coalescing layout     & $0$: none; $1$: row-major; $2$: col-major; $3$: space-filling & 3 \\
$s_3$ & Loop unroll           & $0$: off; $\ell$: factor $2^{\ell}$ & 5 \\
$s_4$ & Reduction strategy    & $0$: scalar; $1$: warp shuffle; $2$: shared-tree; $3$: atomic & 3 \\
$s_5$ & Loop fusion           & $0$: none; $\ell$: fuse $\ell{+}1$ adjacent loops & 3 \\
\bottomrule
\end{tabular}
\end{table}

Rules interact through prerequisites ($\mathcal{E}_{\mathrm{dep}}$),
mutual exclusion ($\mathcal{E}_{\mathrm{excl}}$), and
problem-dependent intensity caps $\bar{M}_i(\mathbf{q})$, yielding the
feasible set
\begin{multline}
\mathcal{A}_0 = \Bigl\{\,
  \mathbf{s} \in \mathcal{A} \mathrel{\Big|}
  \bigwedge_{(i,j)\in \mathcal{E}_{\mathrm{dep}}}
    \bigl(s_j \geq 1 \Rightarrow s_i \geq 1\bigr) \;\land \\
  \bigwedge_{(i,j)\in \mathcal{E}_{\mathrm{excl}}}
    \min(s_i, s_j) = 0 \;\land\;
  \bigwedge_{i=1}^{m} s_i \leq \bar{M}_i(\mathbf{q})
\,\Bigr\}.
\label{eq:feasible}
\end{multline}
A strategy kernel distinguishes ordered from categorical dimensions:
\begin{equation}
\kappa_\mathcal{A}(\mathbf{s},\mathbf{s}')
  = \exp\!\Biggl(-\sum_{i=1}^{m} w_i \, d_i(s_i, s_i')\Biggr),\;\;
d_i(s_i, s_i') = \begin{cases}
|s_i - s_i'|/M_i, & \text{ordered},\\
\mathbb{1}[s_i \neq s_i'], & \text{categorical}.
\end{cases}
\label{eq:skernel}
\end{equation}

\subsection{Performance Space $\mathcal{Y}$}
The reward is the achieved end-to-end speedup, clipped to zero on
failure:
\begin{equation}
y(\mathbf{s}, Q) = \frac{T_{\mathrm{CPU}}(Q)}{T_{\mathrm{GPU}}\bigl(\mathrm{Apply}(Q, \mathbf{s})\bigr)},
\qquad
y = 0 \text{ if compilation fails or output is incorrect}.
\label{eq:reward}
\end{equation}

\subsection{Observation Model}
Each multi-turn experiment yields a tuple
$E_l = (Q_{j(l)}, \mathbf{s}_l^A, \mathbf{s}_l^B, r_l)$, where
$\mathbf{s}^A$ and $\mathbf{s}^B$ are the strategy before and after
an edit. Define the strategy delta
$\Delta \mathbf{s}_l = \mathbf{s}_l^B - \mathbf{s}_l^A$.
We assume a noisy \emph{incremental effect} model:
\begin{equation}
r_l = g\bigl(f(Q_{j(l)}),\; \mathbf{s}_l^A,\; \Delta \mathbf{s}_l\bigr)
      + \epsilon_l, \qquad
\epsilon_l \sim \mathcal{N}(0, \sigma^2).
\label{eq:incremental}
\end{equation}
The three-argument form of $g$ captures that the marginal effect of a
change $\Delta\mathbf{s}$ depends on the current baseline
$\mathbf{s}^A$: \emph{e.g.}, further unroll on top of already-tiled code
behaves differently than unroll on untiled code.

\subsection{Problem Statement}
\label{sec:problem}
\begin{quote}
\textbf{Problem (Strategy Recommendation).}
Given the problem space $\mathcal{P}$ with feature extractor $f$,
the algorithm space $\mathcal{A}$ with feasible subset
$\mathcal{A}_0$, and a history of experiments
$\{E_l\}_{l=1}^{k}$, a new query
$(Q^*, \mathbf{s}^A)$ with feature
$\mathbf{q}^* = f(Q^*)$ (and $\mathbf{s}^A = \mathbf{0}$ when starting
from raw CPU), solve
\begin{equation}
\mathbf{s}^* = \arg\max_{\mathbf{s} \in \mathcal{A}_0}
  \hat{g}\bigl(\mathbf{q}^*,\; \mathbf{s}^A,\; \mathbf{s} - \mathbf{s}^A\bigr),
\label{eq:objective}
\end{equation}
where $\hat{g}$ is the effect-function estimator learned from history.
\end{quote}

This formulation unifies three different views of the problem: (i)~a
\emph{retrieval} view (pick a strategy that worked on a similar
problem); (ii)~a \emph{combinatorial-search} view (greedily optimize
each rule coordinate); and (iii)~a \emph{regression} view (fit
$\hat{g}$ and maximize). Section~\ref{sec:strategy} contrasts the
three.

\section{Benchmark Design}
\label{sec:design}

\subsection{Task Structure}
Each \textsc{AccelEval} task is specified by: \texttt{task.json} (metadata), \texttt{cpu\_reference.c} (CPU implementation), \texttt{prompt\_template.yaml} (natural-language description), \texttt{task\_io.cu} (binary I/O adapter), and \texttt{gen\_data.py} (deterministic input generator).

\subsection{Three-Layer Execution Model}
We decouple correctness from performance by separating the \textbf{harness} (timing and validation), \textbf{task I/O} (binary deserialization), and \textbf{solution} (compute only). Two interface modes govern timing:
\begin{description}
  \item[\texttt{init\_compute}] Large-memory tasks expose separate \texttt{initialize()} and \texttt{compute()} phases; only \texttt{compute()} is timed.
  \item[\texttt{compute\_only}] Small tasks where allocation and transfer are part of the measured region.
\end{description}

\subsection{Correctness Oracle}
The solution's \texttt{output.txt} is compared against an expected output produced by the CPU reference, using either \textbf{exact} string comparison (integer/combinatorial tasks) or \textbf{numerical} comparison via \texttt{numpy.isclose} with task-specific $\mathrm{atol}$ and $\mathrm{rtol}$.

\subsection{Prompting Levels}

\begin{table}[ht]
\centering
\small
\caption{The three prompting regimes. Unless stated otherwise, experiments in Section~\ref{sec:results} use L3.}
\label{tab:levels}
\begin{tabular}{@{}lll@{}}
\toprule
\textbf{Level} & \textbf{Hints given} & \textbf{Purpose} \\
\midrule
L1 & Algorithm background + optimization plan & Measure ceiling with scaffolding \\
L2 & Brief optimization hints & Measure with minimal guidance \\
L3 & None (task description only) & Measure autonomous capability \\
\bottomrule
\end{tabular}
\end{table}

\subsection{Task Distribution}

\begin{table}[ht]
\centering
\small
\caption{Task distribution across algorithmic categories (42 tasks total).}
\label{tab:categories}
\begin{tabular}{@{}lrl@{}}
\toprule
\textbf{Category} & \textbf{\# Tasks} & \textbf{Example} \\
\midrule
Dynamic programming          & 7 & \texttt{network\_rm\_dp}, \texttt{held\_karp\_tsp} \\
Sparse linear algebra        & 4 & \texttt{hpcg\_spmv\_27pt}, \texttt{npb\_cg\_sparse\_solve} \\
Financial computing          & 5 & \texttt{monte\_carlo}, \texttt{black\_scholes} \\
Graph analytics              & 5 & \texttt{gapbs\_pagerank}, \texttt{bellman\_ford} \\
Fluid simulation (SPH)       & 3 & \texttt{sph\_forces}, \texttt{sph\_cell\_index} \\
Spatial distance / clustering& 3 & \texttt{dbscan}, \texttt{hausdorff\_distance} \\
Optimization / LP            & 3 & \texttt{crew\_pairing}, \texttt{pdlp} \\
Stencil simulation           & 2 & \texttt{hotspot\_2d}, \texttt{miniWeather} \\
CFD solver                   & 2 & \texttt{npb\_lu\_ssor\_structured} \\
Other (bioinformatics, stochastic, automata, \ldots) & 8 & --- \\
\midrule
\textbf{Total}               & \textbf{42} & \\
\bottomrule
\end{tabular}
\end{table}

\section{Evaluation Methodology}
\label{sec:eval}

\subsection{Pipeline}
For each \((\text{model}, \text{task}, \text{sample})\) tuple:
\begin{enumerate}
  \item \textbf{Generate}: call the model with the L3 prompt; extract CUDA code from the response.
  \item \textbf{Compile}: \texttt{nvcc -O2 -arch=sm\_89} over \texttt{harness\_gpu.cu + task\_io.cu + solution.cu}.
  \item \textbf{Execute}: run with \texttt{--validate}, producing \texttt{output.txt} and \texttt{timing.json}.
  \item \textbf{Validate}: compare \texttt{output.txt} against the precomputed expected output.
  \item \textbf{Benchmark}: extract mean/std/min/max of CUDA Event timings; report end-to-end and kernel-only speedups over the pinned CPU baseline.
\end{enumerate}

\subsection{Metrics}
\textbf{Compile rate} (fraction compiling), \textbf{Pass rate} (fraction compiling \emph{and} producing correct output), \textbf{Speedup} ($T_{\mathrm{CPU}} / T_{\mathrm{GPU}}$, median and best-of-$k$), and \textbf{$\mathrm{fast}_p$} (fraction with speedup $\geq p\times$, for $p \in \{1, 5, 10\}$).

\subsection{Setup}
All experiments run on a single NVIDIA L40S (\texttt{sm\_89}) with CUDA~12.4 and GCC~11. CPU baselines use a single thread pinned to core~0. Each model generates one sample per task at L3 with temperature~0.7.

\section{Experiments}
\label{sec:results}

\subsection{Overall Results}

Table~\ref{tab:overall} summarizes the six models on L3. Gemini~3.1~Pro leads in pass rate (88.9\%) and $\mathrm{fast}_{10}$~(86.7\%).

\begin{table}[ht]
\centering
\small
\caption{Overall results on \textsc{AccelEval}-42 (L3, single sample). ``Sp.'' denotes end-to-end speedup over the CPU baseline.}
\label{tab:overall}
\begin{tabular}{@{}lrrrrrr@{}}
\toprule
\textbf{Model} & \textbf{Compile} & \textbf{Pass} & \textbf{Median Sp.} & \textbf{$\mathrm{fast}_1$} & \textbf{$\mathrm{fast}_5$} & \textbf{$\mathrm{fast}_{10}$} \\
\midrule
Gemini 3.1 Pro     & 93.3\% & \textbf{88.9\%} & \textbf{316$\times$} & \textbf{91.1\%} & \textbf{88.9\%} & \textbf{86.7\%} \\
Claude Opus 4.6    & \textbf{95.6\%} & 77.8\% & 213$\times$ & 88.9\% & 80.0\% & 80.0\% \\
Qwen 3.6 Plus      & 88.9\% & 64.4\% & 259$\times$ & 80.0\% & 73.3\% & 73.3\% \\
Kimi K2.5          & 78.9\% & 60.5\% & 261$\times$ & 71.1\% & 68.4\% & 68.4\% \\
GLM 5.1            & 82.8\% & 62.1\% & 228$\times$ & 65.5\% & 58.6\% & 58.6\% \\
DeepSeek V3.2      & 55.0\% & 47.5\% & 226$\times$ & 52.5\% & 47.5\% & 45.0\% \\
\bottomrule
\end{tabular}
\end{table}

\subsection{Per-Category Analysis}
Gemini dominates dynamic programming and graph analytics; Claude leads in bioinformatics and selected sparse linear algebra tasks; Qwen and Kimi are competitive on embarrassingly-parallel workloads (Monte~Carlo, stencil).

\subsection{Failure Modes}
We classify failures into three buckets:
\begin{itemize}
  \item \textbf{Compile fail}: missing CUDA include, incorrect kernel signature, invalid template instantiation.
  \item \textbf{Wrong result}: compiled and ran, but output mismatched the oracle---usually race conditions, bad reductions, or boundary errors.
  \item \textbf{Correct but slow}: $<1\times$ speedup---typically excessive synchronization or launch overhead on tiny inputs.
\end{itemize}

\section{Decomposition Analysis: What Makes the Kernel Fast?}
\label{sec:decomp}

Leaderboard metrics say \emph{which} model wins; they do not say \emph{why}. A 300$\times$ speedup could be driven by shared-memory tiling, algorithmic substitution, or trivial precomputation. To attribute performance to concrete techniques we build a two-stage pipeline on top of \textsc{AccelEval}: an \textbf{optimization pattern database} (Section~\ref{sec:pat-db}) and a \textbf{pairwise diff decomposition} over multi-turn agent runs (Section~\ref{sec:pairwise}).

\subsection{Optimization Pattern Database}
\label{sec:pat-db}

\begin{figure}[ht]
\centering
% \includegraphics[width=\linewidth]{figs/pattern_db.pdf}
\caption{Three-stage pipeline for the pattern database. Auto-detect uses static CUDA analysis (\texttt{\_\_shared\_\_}, \texttt{\_\_constant\_\_}, \texttt{template<>}, \texttt{\#pragma unroll}, \texttt{atomicAdd}, \texttt{\_\_shfl\_sync}, \ldots); an LLM analyzer reads the CPU reference and generated CUDA to summarize detected patterns and propose new candidates; the knowledge base accumulates evidence across runs/models, with auto-promotion once a candidate appears in sufficiently many passing samples.}
\label{fig:patdb}
\end{figure}

Each sample that compiles and passes validation is parsed by two complementary analyzers:
(i) a \emph{static detector} scans for canonical CUDA idioms (shared-memory usage, warp intrinsics, template dispatch, atomic ops, restricted pointers, \texttt{\#pragma unroll}, etc.); and
(ii) an \emph{LLM analyzer} (Gemini~3.1~Pro, budget-bounded) reads both the CPU reference and the generated CUDA and produces a natural-language pattern list, including novel candidates not in the registry.
Detected patterns are written into a pattern database with evidence counts and speedup annotations. Candidates are auto-promoted to first-class patterns once they accumulate evidence across multiple models and tasks.

\begin{table}[ht]
\centering
\small
\caption{Example patterns in the database (selected).}
\label{tab:patterns}
\begin{tabular}{@{}llrl@{}}
\toprule
\textbf{ID} & \textbf{Pattern} & \textbf{Evidence} & \textbf{Observed lift} \\
\midrule
PAT-001 & Shared-memory tiling                & 42 & $12.3\times$ avg lift \\
PAT-005 & Bitmask state representation        & 8  & $3{,}186\times$ (NFA) \\
PAT-008 & Template kernel dispatch            & 15 & $8.7\times$ avg lift \\
PAT-012 & Precomputation ($\varepsilon$-closure fusion) & 6  & eliminates BFS phase \\
\bottomrule
\end{tabular}
\end{table}

Across the 40 tasks that produced a passing solution, our agent discovered \textbf{33 distinct optimization strategies}, with the most frequent being \emph{persistent device-memory allocation} (35 tasks), \emph{precomputation via init kernel} (32 tasks), and \emph{memory coalescing} (31 tasks). The long tail contains task-specific techniques (\emph{e.g.}, $\varepsilon$-closure fusion for regex, grid-stride loops for reductions) that appear in only one or two tasks. A strategy-intensity matrix (tasks~$\times$~patterns) reveals which categories are ``saturated'' by a few common patterns and which require bespoke strategies.

\subsection{Pairwise Diff Attribution}
\label{sec:pairwise}

\textsc{AccelEval} supports a \emph{multi-turn agent} mode in which the model is given feedback (compilation errors, correctness mismatches, slow kernel profile) and produces successive versions $V_0, V_1, \ldots, V_k$ of the same solution. Consecutive versions form natural A/B experiments: the same task, same oracle, and two code snapshots differing by a small edit. We attribute speedup change to the concrete patterns that changed between them.

Formally, for each consecutive pair $(V_a, V_b)$ we record a \texttt{DiffRecord}:
\begin{itemize}
  \item \texttt{pattern\_changes[]}: patterns added/removed/modified between $V_a$ and $V_b$;
  \item \texttt{causal\_chains[]}: ordered sequences inferred across multiple diffs (\emph{e.g.}, shared-memory tiling enables subsequent template specialization);
  \item \texttt{new\_candidates[]}: previously unseen techniques flagged by the LLM analyzer;
  \item \texttt{speedup\_ratio} $= T(V_a) / T(V_b)$.
\end{itemize}

Aggregating \texttt{DiffRecord}s across all multi-turn traces yields a \emph{pattern$\to$speedup attribution table}. In one representative trace ($V_0 = 5.2\times$, $V_4 = 48.7\times$), the decomposition attributes $4.1\times$ of the $9.4\times$ end-to-end improvement to shared-memory tiling, $2.3\times$ to template dispatch, and the remainder to bounds-check elimination.

\paragraph{Why pairwise (not absolute).}
Absolute correlations between ``pattern present'' and ``speedup'' are confounded: fast kernels often bundle many patterns. Pairwise diffs isolate the marginal contribution of each edit, approximating the counterfactual ``what if only this pattern changed?'' without requiring manual ablations.

\paragraph{Limitations.}
Attribution is only as reliable as the granularity of the agent's edits. Large, multi-pattern diffs provide weaker causal claims than small ones. Our current pipeline reports both per-pattern lifts and bundle-level lifts for transparency.

\section{Strategy Recommendation: Three Baselines}
\label{sec:strategy}

Using the formulation of Section~\ref{sec:formulation}, we benchmark
three baselines of increasing complexity. All three share the same
feature extractor $f$ and feasible set $\mathcal{A}_0$; they differ
only in how they estimate the effect function $\hat{g}$ in
Equation~\eqref{eq:objective}.

\subsection{Baseline 1: Nearest-Neighbor Transfer}
\label{sec:bl1}
Retrieve the most similar problem in the history and reuse its
best-observed strategy:
\begin{equation}
j^* = \arg\min_{j}\;\bigl\lVert f(Q^*) - f(Q_j)\bigr\rVert_2,
\qquad
\mathbf{s}^* = \arg\max_{\substack{\mathbf{s}:\;\exists\,l \text{ s.t. } j(l)=j^*,\; \mathbf{s}_l^B = \mathbf{s}}} r_l.
\label{eq:nn}
\end{equation}
We also report the $k$-NN variant with distance-weighted voting. This
baseline learns no model; it relies purely on the locality assumption
in feature space.

\subsection{Baseline 2: Greedy Coordinate Search}
\label{sec:bl2}
Starting from $\mathbf{s} = \mathbf{s}^A$, sweep one dimension at a
time. In round $i$, hold all other coordinates fixed and choose the
intensity that maximizes observed speedup:
\begin{equation}
s_i^* = \arg\max_{s_i \in \{0,\ldots,M_i\}}\,
  y\bigl(\mathbf{s}[s_i \!\to\! s_i],\; Q^*\bigr),
  \qquad i = 1, \ldots, m.
\label{eq:greedy}
\end{equation}
Total evaluations scale as $\sum_{i=1}^{m}(M_i + 1)$, compared to the
exhaustive $\prod_i(M_i+1)$. Greedy search ignores cross-dimension
interactions but costs only linear evaluations.

\subsection{Baseline 3: Linear Regression}
\label{sec:bl3}
Fit the effect function as a linear model in features and strategy:
\begin{equation}
\hat{g}(\mathbf{q}, \mathbf{s}^A, \Delta\mathbf{s})
  = \mathbf{w}_q^{\top}\mathbf{q}
  + \mathbf{w}_s^{\top}\mathbf{s}^A
  + \mathbf{w}_\Delta^{\top}\Delta\mathbf{s} + b,
\label{eq:linreg}
\end{equation}
learned by least squares on $\{(\mathbf{x}_l, r_l)\}_{l=1}^{k}$.
Substituting into Equation~\eqref{eq:objective} gives a separable
optimum per coordinate (or an integer-linear program when
$\mathcal{A}_0$ is binding). Linear regression cannot capture
interactions but provides an interpretable, principled lower bound.

\subsection{Experimental Protocol}

\paragraph{Data split.} We hold out $20\%$ of $\mathcal{P}$ as
unseen queries (``test tasks'') and train all three baselines on the
remaining $80\%$, using the multi-turn agent traces from
Section~\ref{sec:decomp} as the history $\{E_l\}$.

\paragraph{Metrics.}
\begin{itemize}
  \item \textbf{Hit@$K$} --- fraction of test queries whose predicted
  $\mathbf{s}^*$ (top-$K$) matches the best-observed strategy in the
  held-out trace.
  \item \textbf{Speedup gap} --- $\Delta = y(\mathbf{s}^{\mathrm{oracle}})
  - y(\mathbf{s}^*)$, averaged over test queries.
  \item \textbf{Evaluation budget} --- number of compile-and-run trials
  needed to reach within $10\%$ of the oracle speedup.
\end{itemize}

\subsection{Results}
\label{sec:strategy-results}

\begin{table}[ht]
\centering
\small
\caption{Strategy-recommendation baselines on \textsc{AccelEval}-42.
\emph{(Results to be filled.)} $\mathbf{s}^A = \mathbf{0}$ (raw CPU);
history is drawn from the multi-turn traces of
Section~\ref{sec:decomp}.}
\label{tab:strategy}
\begin{tabular}{@{}lrrr@{}}
\toprule
\textbf{Method} & \textbf{Hit@1} & \textbf{Avg.\ speedup gap} & \textbf{\# trials to 0.9$\times$oracle} \\
\midrule
Baseline 1: 1-NN Transfer       & --- & --- & --- \\
Baseline 1$^\prime$: 5-NN Voting & --- & --- & --- \\
Baseline 2: Greedy Coordinate    & --- & --- & --- \\
Baseline 3: Linear Regression    & --- & --- & --- \\
\midrule
LLM (Gemini 3.1 Pro, L3)          & --- & --- & --- \\
LLM + history (retrieval-augmented) & --- & --- & --- \\
\bottomrule
\end{tabular}
\end{table}

\begin{figure}[ht]
\centering
% \includegraphics[width=0.7\linewidth]{figs/strategy_budget.pdf}
\caption{\emph{(To be filled.)} Speedup attained versus
compile-and-run budget for the three baselines and the LLM.
Expected trend: Baseline~1 saturates early at its retrieval ceiling;
Baseline~2 improves monotonically but plateaus before the oracle;
Baseline~3 is competitive at very low budgets but asymptotes below
the interaction-aware methods.}
\label{fig:strategy-budget}
\end{figure}

\paragraph{What we expect.}
We anticipate: (i)~NN-transfer will dominate on tasks with close
neighbors in $\mathcal{F}$ (graph analytics, dense linear algebra),
but collapse on rare categories (regex, combinatorial optimization);
(ii)~greedy coordinate search will outperform NN-transfer given
a modest budget, but miss interactions such as tiling$\times$unroll;
(iii)~linear regression will be the weakest on raw speedup but most
sample-efficient; (iv)~strong LLMs will roughly match or exceed the
classical baselines on tasks close to pretraining data, and fall
sharply on bespoke categories.

\section{Strategy Transfer: Discovery vs.\ Implementation}
\label{sec:transfer}

The L3 results in Section~\ref{sec:results} conflate two skills:
\emph{knowing what to optimize} (strategy discovery) and
\emph{knowing how to code it} (implementation). A model could
underperform because it cannot identify useful patterns, or because
it identifies them but writes buggy/slow CUDA. We design a two-stage
ablation to tell these apart.

\paragraph{Stage~1 (baseline, from Section~\ref{sec:results}).}
Each model generates a solution under L3. Let $M^{*}(Q)$ denote the
best model on task $Q$ (highest correct speedup).

\paragraph{Stage~2 (strategy transfer).}
For each non-winning model $M$ on task $Q$, we take $M^{*}(Q)$'s
solution, extract its optimization strategy as a natural-language
recipe (ordered list of patterns + rationale), and re-prompt $M$ with
this recipe appended to the L3 prompt. $M$ must then produce its own
implementation guided by the recipe, without seeing $M^{*}$'s code.

\paragraph{Diagnosis.} Let $y^{\mathrm{S1}}_{M}$ and
$y^{\mathrm{S2}}_{M}$ be the speedups in Stage~1 and Stage~2. Define
the \emph{replication ratio}
\begin{equation}
\mathrm{RR}(M, Q) = \frac{y^{\mathrm{S2}}_{M}(Q)}{y(M^*)(Q)}
  \in [0, 1].
\label{eq:rr}
\end{equation}
We read the outcome per task:
\begin{itemize}
  \item $\mathrm{RR} \geq \tau$ (large, close to~1): the Stage~1 gap
        was \textbf{strategy discovery} --- given the recipe, $M$
        can execute.
  \item $\mathrm{RR} < \tau$ (small): the gap is \textbf{implementation
        skill} --- even told what to do, $M$ cannot write the code.
\end{itemize}
We use $\tau = 0.7$ (within $30\%$ of the leader) as the default.

\paragraph{Why this matters.}
The two diagnoses call for different remedies. A discovery gap points
to retrieval augmentation (Section~\ref{sec:strategy}) or agentic
reasoning. An implementation gap is harder: it points to limitations
in the base model's CUDA fluency, which supervised fine-tuning or
longer context with canonical examples may fix.

\paragraph{Experimental protocol.}
\begin{enumerate}
  \item Select the $T$ tasks where Gemini~3.1~Pro (the Stage~1
        leader) achieves $\geq 10\times$ speedup and no other model is
        within $2\times$. ($T \approx 15$ in our current data.)
  \item For each competitor $M \in \{$Claude, Qwen, Kimi, GLM,
        DeepSeek$\}$, run the transfer protocol.
  \item Report per-model, per-task RR and the fraction of tasks
        classified as \emph{discovery} vs.\ \emph{implementation}.
\end{enumerate}

\begin{table}[ht]
\centering
\small
\caption{Strategy-transfer diagnosis per competitor model.
\emph{(Results to be filled.)} For each model we report the mean
replication ratio $\mathrm{RR}$, the fraction of tasks with
$\mathrm{RR} \geq 0.7$ (``discovery gap''), and the fraction with
$\mathrm{RR} < 0.7$ (``implementation gap'').}
\label{tab:transfer}
\begin{tabular}{@{}lrrrr@{}}
\toprule
\textbf{Competitor} & \textbf{\# tasks} & \textbf{Mean RR} & \textbf{Discovery\%} & \textbf{Impl.\%} \\
\midrule
Claude Opus 4.6 & --- & --- & --- & --- \\
Qwen 3.6 Plus    & --- & --- & --- & --- \\
Kimi K2.5        & --- & --- & --- & --- \\
GLM 5.1          & --- & --- & --- & --- \\
DeepSeek V3.2    & --- & --- & --- & --- \\
\bottomrule
\end{tabular}
\end{table}

\begin{figure}[ht]
\centering
% \includegraphics[width=0.55\linewidth]{figs/transfer_diagnosis.pdf}
\caption{\emph{(To be filled.)} Scatter of $y^{\mathrm{S1}}$ (x-axis)
vs.\ $y^{\mathrm{S2}}$ (y-axis) per task. Points on the $y = y^{\mathrm{S2}}_{\mathrm{leader}}$ line indicate a fully-closed discovery gap;
points clustered near the $y = y^{\mathrm{S1}}$ line indicate a
stubborn implementation gap.}
\label{fig:transfer}
\end{figure}

\section{Discussion}
\label{sec:discussion}

\paragraph{Where LLMs excel.} Models consistently produce strong CUDA code for massively parallel, regular workloads such as Monte~Carlo simulation and dense element-wise kernels, often exceeding $1000\times$ speedup.

\paragraph{Where they fail.} Tasks requiring irregular memory access patterns (sparse solvers, graph irregularity) or cross-iteration dependencies (multigrid V-cycles, Gauss-Seidel sweeps) remain fragile. Three tasks (\texttt{nbnxm\_forces}, \texttt{npb\_mg\_vcycle}, \texttt{rodinia\_nw\_alignment}) were failed by \emph{all} six models.

\paragraph{Infrastructure matters.} Routing Chinese models through OpenRouter introduces significant instability: $\sim$18\% of GLM requests returned empty responses in our runs. We added empty-response retries and throughput-based routing to partially mitigate this.

\section{Limitations}
\label{sec:limitations}

\begin{itemize}
  \item \textbf{Single sample per task.} We report one generation per $(\text{model},\text{task})$ pair; $k$-sample studies would tighten pass-rate confidence intervals.
  \item \textbf{Single GPU architecture.} All measurements on \texttt{sm\_89}; generalization to H100 / A100 / consumer cards is unstudied.
  \item \textbf{L3 only.} We have not yet scaled the L1/L2 evaluation across all models.
\end{itemize}

\section{Conclusion}
\label{sec:conclusion}

\textsc{AccelEval} provides a rigorous, reproducible platform for evaluating LLM-generated GPU code on realistic workloads. Our first large-scale study shows a clear capability gap between frontier closed-source models and open-weight alternatives, concentrated in correctness on irregular HPC kernels. We release the benchmark, the harness, and the evaluation scripts to support community benchmarking.

\section*{Broader Impact}

\textsc{AccelEval} aims to make LLM-assisted GPU programming more rigorous and reproducible. Positive impact: lowering the barrier to entry for GPU acceleration in scientific computing; enabling fair cross-model evaluation; surfacing failure modes that model developers can target. Potential concerns: benchmark saturation incentivizing overfitting; uneven access to GPU infrastructure may advantage models that produce simpler kernels. We release all task definitions, harnesses, and evaluation logs to enable independent verification.



\begin{ack}
We thank \emph{(TBD)}. No external funding to declare. No competing interests.
\end{ack}

\section*{References}

\medskip
{\small

\bibitem[Ouyang et al.(2025)]{kernelbench}
A.\ Ouyang, S.\ Guo, S.\ Arora, A.\ Zhang, M.\ Hu, C.\ R\'e, and A.\ Mirhoseini.
(2025) KernelBench: Can LLMs write efficient GPU kernels?
{\it arXiv preprint arXiv:2502.10517}.

\bibitem[Chen et al.(2021)]{humaneval}
M.\ Chen et al.\ (2021) Evaluating large language models trained on code.
{\it arXiv preprint arXiv:2107.03374}.

\bibitem[Austin et al.(2021)]{mbpp}
J.\ Austin et al.\ (2021) Program synthesis with large language models.
{\it arXiv preprint arXiv:2108.07732}.

\bibitem[Jimenez et al.(2024)]{swebench}
C.\ E.\ Jimenez et al.\ (2024) SWE-bench: Can language models resolve
real-world GitHub issues?
{\it ICLR}.

\bibitem[Wen et al.(2022)]{babeltower}
Y.\ Wen et al.\ (2022) BabelTower: Learning to auto-parallelized program
translation.
{\it ICML}.

\bibitem[Tehranijamsaz et al.(2024)]{coderosetta}
A.\ Tehranijamsaz et al.\ (2024) CodeRosetta: Pushing the boundaries of
unsupervised code translation for parallel programming.
{\it NeurIPS}.

\bibitem[Nugteren \& Corporaal(2014)]{bones}
C.\ Nugteren and H.\ Corporaal.\ (2014) Bones: An automatic skeleton-based
C-to-CUDA compiler for GPUs.
{\it ACM TACO}, 11(4).

\bibitem[Amini et al.(2012)]{par4all}
M.\ Amini et al.\ (2012) Par4All: From convex array regions to heterogeneous
computing.
{\it IMPACT}.
}

%%%%%%%%%%%%%%%%%%%%%%%%%%%%%%%%%%%%%%%%%%%%%%%%%%%%%%%%%%%%

\appendix

\section{Complete Task List}
\label{app:tasklist}

\emph{(42 tasks with category, difficulty, and brief description.
Auto-generated from \texttt{tasks/*/task.json}.)}

\section{Example Prompt (L3)}
\label{app:prompt}

\begin{lstlisting}[style=cuda,caption={Canonical L3 prompt structure.}]
# Task: <name>

<task description>

## Interface
<C function signatures>

## Input sizes
<auto-generated size table>

## Output contract
<expected output format>

## CPU Reference
<cpu_reference.c content>

Please produce a single self-contained CUDA file that implements
<function names>. Your code will be compiled and linked against
the provided harness.
\end{lstlisting}

%%%%%%%%%%%%%%%%%%%%%%%%%%%%%%%%%%%%%%%%%%%%%%%%%%%%%%%%%%%%

\newpage
\input{checklist.tex}


\end{document}